%\section[Ottica Geometrica: Riflessione e Specchi]
    \begin{frame}{L'Ottica Geometrica}
    \begin{itemize}
        \item Regime valido quando $\lambda \ll$ dimensioni degli ostacoli
        \item La luce viene modellata come raggio rettilineo
        \item Si studiano: riflessione (specchi) e rifrazione (lenti, prismi)
        \item Convenzione dei segni (da applicare sempre):
        \begin{itemize}
            \item Distanze reali (davanti alla superficie): positive
            \item Distanze virtuali (dietro la superficie): negative
        \end{itemize}
        
    \end{itemize}
\end{frame}

%========================================================

\begin{frame}{La Legge di Riflessione}
    \begin{itemize}
        \item La luce che colpisce una superficie liscia viene riflessa
        \item Legge di riflessione:
        $$\theta_i = \theta_r$$
        \item Entrambi gli angoli si misurano rispetto alla normale alla superficie
        \item Raggio incidente, normale e raggio riflesso giacciono sullo stesso piano
        \item Riflessione speculare (superficie liscia) $\neq$ Riflessione diffusa (superficie rugosa)
    \end{itemize}
\end{frame}

%========================================================

\begin{frame}{\href{https://bhuiyanek.github.io/Fisica/riflessione.html?mode=piano}{Specchio Piano}}
    \begin{itemize}
        \item L'immagine è virtuale: i raggi non convergono fisicamente dietro lo specchio
        \item L'immagine è diritta e delle stesse dimensioni dell'oggetto
        \item Distanza immagine = distanza oggetto (simmetria rispetto alla superficie)
        \item Lateralmente invertita (destra $\leftrightarrow$ sinistra) 
    \end{itemize}
\end{frame}


%========================================================

\begin{frame}{Specchi Sferici: Grandezze Fondamentali}
    \begin{itemize}
        \item Centro di curvatura $C$: centro della sfera di cui lo specchio è parte
        \item Raggio di curvatura $R$: distanza da $C$ alla superficie
        \item Asse ottico: retta passante per $C$ perpendicolare allo specchio
        \item Fuoco $F$: punto in cui convergono i raggi paralleli all'asse ottico
        \item Distanza focale: $f = \frac{R}{2}$
        \item Specchio concavo: $f > 0$ — Specchio convesso: $f < 0$
    \end{itemize}
\end{frame}

%========================================================

\begin{frame}{\href{https://bhuiyanek.github.io/Fisica/riflessione.html?mode=concavo}{Specchio Concavo}}

    \begin{itemize}
        \item Superficie riflettente rivolta verso il \textbf{centro di curvatura}
        \item I raggi paralleli all’asse ottico convergono nel \textbf{fuoco reale $F$}
        \item La natura dell’immagine dipende dalla posizione dell’oggetto:
    \end{itemize}
    
    \medskip
    \centering
    
    \renewcommand{\arraystretch}{1.3}
    
    \begin{tabular}{|l|l|}
    \hline
    \rowcolor{maincolor!20}
    \textbf{Posizione oggetto} & \textbf{Immagine} \\
    \hline
    Oltre $C$ & Reale, capovolta, rimpicciolita \\
    \hline
    In $C$ & Reale, capovolta, stessa dimensione \\
    \hline
    Tra $C$ e $F$ & Reale, capovolta, ingrandita \\
    \hline
    In $F$ & Immagine all'infinito \\
    \hline
    Entro $F$ & Virtuale, diritta, ingrandita \\
    \hline
    \end{tabular}
    
    \medskip
    
    Applicazioni: telescopi riflettori, specchi per trucco, concentratori solari.

\end{frame}

%========================================================

\begin{frame}{\href{https://bhuiyanek.github.io/Fisica/riflessione.html?mode=convesso}{Specchio Convesso}}
    \begin{itemize}
        \item Superficie riflettente rivolta verso l'esterno
        \item I raggi paralleli divergono dopo la riflessione
        \item Il fuoco è virtuale (dietro lo specchio), $f < 0$
        \item L'immagine è sempre: virtuale, diritta, rimpicciolita
        \item Vantaggio: campo visivo molto ampio
        \item Applicazioni: retrovisori, specchi stradali, specchi di sorveglianza
    \end{itemize}
\end{frame}

%========================================================
\begin{frame}{Equazione dei punti coniugati}

    \textbf{Relazione fondamentale (specchi sferici):}
    
    \medskip
    
    \[
    \frac{1}{p} + \frac{1}{p'} = \frac{1}{f} = \frac{2}{R}
    \]
    
    \medskip
    
    \textbf{Convenzioni dei segni:}
    
    \begin{itemize}
        \item $p > 0$: oggetto reale (davanti allo specchio)
        \item $p' > 0$: immagine reale (davanti allo specchio)
        \item $p' < 0$: immagine virtuale (dietro lo specchio)
    \end{itemize}

\end{frame}

\begin{frame}{Ingrandimento lineare}

    \textbf{Definizione:}
    
    \medskip
    
    \[
    G = -\frac{p'}{p} = \frac{y'}{y}
    \]
    
    \medskip
    
    \textbf{Interpretazione fisica:}
    
    \begin{itemize}
        \item $|G| > 1$ → immagine ingrandita
        \item $|G| < 1$ → immagine rimpicciolita
        \item $G > 0$ → immagine diritta (virtuale)
        \item $G < 0$ → immagine capovolta (reale)
    \end{itemize}

\end{frame}